%% appendix1.tex — Appendix A: Ad Hoc Interface Mode in Linux Systems

\chapter{Ad Hoc Interface Mode in Linux Systems}
\label{app:adhoc}

This appendix documents how the wireless interface of each node is placed in
ad hoc (IBSS) mode, how the network used throughout the evaluation is set up, and
how the kernel routing tables that carry the mesh traffic are configured. It
collects, in one place, the steps that the deployment scripts of this work
(\texttt{adhoc-start.sh} and the node software) perform automatically, so that
the setup can be reproduced on a fresh system.

%% -----------------------------------------------------------------------
\section{Active Nodes}
\label{sec:app-nodes}

Table~\ref{tab:app-nodes} lists the active nodes of the testbed, with their
role, hardware model, the physical address assigned to the wireless interface
(\texttt{wlan0}, in the \texttt{172.20.10.0/28} ad hoc plane), the mesh address
assigned to the TUN interface (in the \texttt{10.0.0.0/24} overlay), and the MAC
address of the wireless interface.

\begin{table}[h]
\centering
\caption{Network information of the active nodes}
\label{tab:app-nodes}
\begin{tabular}{lllll}
\toprule
\textbf{Node} & \textbf{Model} & \textbf{Physical IP} & \textbf{Mesh IP} & \textbf{MAC Address} \\
\midrule
N1 (Robot)        & AlphaBot2-Pi / RPi~5 & 172.20.10.1 & 10.0.0.1 & \texttt{xx:xx:xx:xx:xx:xx} \\
N2 (Relay)        & Raspberry Pi~5       & 172.20.10.2 & 10.0.0.2 & \texttt{xx:xx:xx:xx:xx:xx} \\
N3 (Base Station) & Laptop (Ubuntu 22.04)& 172.20.10.3 & 10.0.0.3 & \texttt{xx:xx:xx:xx:xx:xx} \\
\bottomrule
\end{tabular}
\end{table}

%% -----------------------------------------------------------------------
\section{Configuration of Ad Hoc Interfaces}
\label{sec:app-config}

The first step is to confirm that the wireless device supports ad hoc mode.
Assuming the interface is named \texttt{wlan0}, this is checked with:

\begin{lstlisting}[language=bash]
iw dev wlan0 info
\end{lstlisting}

The output describes the \texttt{wlan0} interface and, under the entry
``Supported interface modes'', lists ``IBSS'' (equivalently ``Ad-Hoc'') when ad
hoc mode is available.

Before the mode can be changed, no service must be holding control of the
interface, because an active network manager will revert manual changes. If the
only manager is NetworkManager, the mode can be changed once it is stopped. If a
service such as \texttt{dhcpcd} or \texttt{wpa\_supplicant} is active, it must be
stopped or disabled first, as it would otherwise reclaim the interface and
restrict manual configuration. The state of \texttt{dhcpcd}, for example, can be
checked with \texttt{sudo systemctl status dhcpcd}. When \texttt{wpa\_supplicant}
is in use, adding the line \texttt{ap\_scan=2} to
\texttt{/etc/wpa\_supplicant/wpa\_supplicant.conf} lets the system keep scanning
for networks even when none is configured.

\subsection{Configuration with the dhcpcd Service}
\label{subsec:app-dhcpcd}

If \texttt{dhcpcd} (or a similar service) is active, there are two options:
either disable the service and reboot, or install NetworkManager and start it
only when needed.

\begin{itemize}
  \item \textbf{Disable dhcpcd.} Run \texttt{sudo systemctl disable dhcpcd} and
  reboot. The interface mode can then be changed freely.
  \item \textbf{Install NetworkManager.} Install NetworkManager and start it on
  demand (for example from a start-up script), with \texttt{sudo systemctl start
  NetworkManager}. This keeps an infrastructure connection available for SSH
  access while still allowing ad hoc mode to be enabled on the same or another
  interface when required. Code Block~\ref{lst:app-nm} installs the service and
  its graphical front end.
\end{itemize}

\begin{lstlisting}[language=bash,caption={Installation of the NetworkManager service},label=lst:app-nm]
sudo apt update
sudo apt install network-manager network-manager-gnome \
     openvpn network-manager-openvpn network-manager-openvpn-gnome -y
\end{lstlisting}

\subsection{Changing the Interface Mode to Ad Hoc}
\label{subsec:app-mode}

Once no service is managing the interface restrictively, the mode can be changed
to ad hoc. This work uses the ESSID \texttt{manet-mesh} on channel~6, with each
node taking the physical address \texttt{172.20.10.X}, where \texttt{X} is the
node identifier.

\paragraph{Manually.} Code Block~\ref{lst:app-manual} shows the commands used by
the deployment script \texttt{adhoc-start.sh} to bring the interface up in ad hoc
mode with the network used in this dissertation.

\begin{lstlisting}[language=bash,caption={Setting up the ad hoc network manually},label=lst:app-manual]
systemctl stop NetworkManager
ip link set wlan0 down
iwconfig wlan0 mode ad-hoc
iwconfig wlan0 essid manet-mesh
iwconfig wlan0 channel 6
ip link set wlan0 up
ip addr add 172.20.10.X/28 dev wlan0
\end{lstlisting}

\paragraph{Automatically.} On the nodes, these commands run unattended at boot
from a \texttt{systemd} unit (\texttt{adhoc.service}) scheduled before the mesh
software and before \texttt{wpa\_supplicant} can interfere. Alternatively, where
no service enforces infrastructure mode, the interface can be brought up in ad
hoc mode directly by the system's network configuration. Code
Block~\ref{lst:app-interfaces} shows the equivalent entry in
\texttt{/etc/network/interfaces}.

\begin{lstlisting}[language=bash,caption={Static ad hoc configuration in the interfaces file (X is the node ID)},label=lst:app-interfaces]
auto wlan0
iface wlan0 inet static
    address 172.20.10.X
    netmask 255.255.255.240
    wireless-mode ad-hoc
    wireless-essid manet-mesh
    wireless-channel 6
\end{lstlisting}

Even when another service overrides this file, it is still worth editing it with
the content above, because the configuration lets the interface keep a static IP
address of its own.

%% -----------------------------------------------------------------------
\section{Routing Tables}
\label{sec:app-routing}

On top of the ad hoc interface, the node software installs the host routes that
carry the mesh overlay. It uses two kernel routing tables. The \emph{main} table
holds, for every reachable mesh destination, a host route (\texttt{/32}) toward
the current next hop, so that ordinary traffic to a \texttt{10.0.0.X} address is
forwarded along the path chosen by the spanning tree. A separate \emph{relay}
table (table~200) holds the routes used when this node forwards traffic on behalf
of others, keeping transit forwarding separate from locally originated traffic.
Both tables are populated directly through the Netlink interface rather than by
spawning \texttt{ip} processes, which keeps each route update at roughly
$50\,\mu$s instead of the few milliseconds an \texttt{ip route}
fork/exec would cost.

The routes installed are equivalent to the following, where \texttt{<dst>} is a
mesh destination and \texttt{<next\_hop>} the chosen relay:

\begin{lstlisting}[language=bash,caption={Equivalent routing-table entries (installed via Netlink)},label=lst:app-routes]
# main table: host route to each mesh destination via its next hop
ip route add 10.0.0.<dst>/32 via 10.0.0.<next_hop> dev tun

# relay table (200): forwarding entries for transit traffic
ip route add 10.0.0.<dst>/32 via 10.0.0.<next_hop> dev tun table 200

# the relay table is cleared and rebuilt whenever the topology changes
ip route flush table 200
\end{lstlisting}

When the topology changes, the relay table is flushed and rebuilt from the new
spanning tree, while the main table entries are updated in place, so that the
forwarding state always reflects the current routes.
